
\section{Conclusion}\label{sec:conclusion}

We studied speculative decoding as a serving decision in which using and improving a drafter must be weighed against their costs under the current workload. Our measurements show that better acceptance does not necessarily translate into higher serving throughput: the profitability of speculation changes with concurrency and target architecture, while online adaptation consumes resources that would otherwise serve requests. To address this trade-off, we developed a workload-aware controller that combines a calibrated profitability frontier, drafter learning estimates, and training opportunity costs to decide when to speculate, how much to train, and when to stop. Across three workloads spanning dense and MoE targets, the controller stays within \nv{full.worstshort}\% of the best measured policy in every scenario, without knowing the optimal training stopping point in advance. The ablations reveal that different workloads benefit from different decisions: disabling unprofitable speculation is crucial for Chinese traffic, controlling training interference matters for dense code, and continued adaptation benefits the MoE target. Mixed-traffic experiments further show that a drafter capable of crossing the profitability boundary may still cost more to train than it returns, while the SGLang deployment demonstrates that correct control decisions remain subject to engine-specific execution overheads. 
The goal is not to keep improving the drafter, but to improve serving performance when adaptation pays for itself.

